\documentclass[11pt]{article}
\usepackage[margin=1in]{geometry}
\usepackage{amsmath,amssymb,booktabs,hyperref}
\title{Replication: \emph{Nearly Perfect Spin Polarization of Noncollinear Antiferromagnets} (Mn$_3$GaN)\\ \large Gurung, Elekhtiar, Luo, Shao \& Tsymbal (2023)}
\author{Textures-100 automated replication (Hermes subagent)}
\date{\today}
\begin{document}
\maketitle

\begin{abstract}
We build a from-scratch tight-binding surrogate of the paper's illustrative
noncollinear-antiferromagnet model and verify the headline claim: a noncollinear
AFM on a kagome lattice supports conduction channels with \textbf{nearly 100\%}
effective momentum-dependent spin polarization. Our model reproduces (i) the
six spin-split bands without spin--orbit coupling and (ii) a maximum effective
spin polarization $p_{k\parallel}^{\max}=0.99997$, i.e.\ essentially 100\%,
over a substantial region of the $(k_y,E_F)$ plane. \textbf{Verdict: REPLICATED}
for the tight-binding headline; the full DFT Mn$_3$GaN Fermi surface and the
ETMR $\sim10^4\%$ transmission calculation are out of scope (DFT skipped by design).
\end{abstract}

\section{Paper claim}
The authors define an effective momentum-dependent spin polarization for the
conduction channels of a metal electrode. With transport along $z$ and conserved
transverse momentum $\mathbf{k}_\parallel$, the net channel spin is
$\mathbf{s}_{\mathbf{k}_\parallel}=\sum_n \mathbf{s}_{n\mathbf{k}_\parallel}$ and
\begin{equation}
\mathbf{p}_{\mathbf{k}_\parallel}=\frac{\mathbf{s}_{\mathbf{k}_\parallel}}
{\sum_n |\mathbf{s}_{n\mathbf{k}_\parallel}|},\qquad
p_{\mathbf{k}_\parallel}\equiv|\mathbf{p}_{\mathbf{k}_\parallel}|.
\tag{Eq.~2}
\end{equation}
$p_{\mathbf{k}_\parallel}=100\%$ when all channel spins at $\mathbf{k}_\parallel$
are parallel, or when only one channel is present. The headline: noncollinear
antiferromagnet Mn$_3$GaN (antiperovskite, $\Gamma_{5g}$ 120$^\circ$ spin texture)
exhibits nearly 100\% spin-polarized states over a broad area of the Fermi surface,
enabling ETMR as large as $10^4\%$ through SrTiO$_3$.

\section{Model built (from scratch)}
We reproduce the paper's Fig.~1 illustrative model: a 2D \textbf{kagome} lattice
(3 sublattices $\times$ 2 spin $=6$ bands), one orbital per atom, spin-independent
nearest-neighbor hopping $t$, and on-site exchange splitting $\Delta$ aligned to a
120$^\circ$ noncollinear (in-plane, $\Gamma_{5g}$-like) AFM moment configuration:
\begin{equation}
H(\mathbf{k}) = H_{\text{hop}}(\mathbf{k})\otimes I_2
+ \frac{\Delta}{2}\sum_{i=A,B,C} P_i\otimes(\mathbf{m}_i\cdot\boldsymbol{\sigma}),
\quad \mathbf{m}_i \in \{90^\circ,210^\circ,330^\circ\}.
\end{equation}
The 120$^\circ$ texture breaks $\hat P\hat T$ and $\hat T\hat t$, giving spin-split
bands \emph{without} SOC (the paper's central symmetry requirement). We set
$\Delta/t=1.5$ (paper value). For transport along $x$, $k_\parallel=k_y$; for each
$(k_y,E_F)$ we collect all Fermi crossings $E_n(k_x)=E_F$, take their spin
expectations $\mathbf{s}_n=\langle\psi_n|\boldsymbol{\sigma}/2|\psi_n\rangle$, and
evaluate Eq.~2.

\section{Results}
\begin{table}[h]\centering
\begin{tabular}{lll}
\toprule
Quantity & This work & Paper \\
\midrule
Number of bands & 6 & 6 \\
Spin-split without SOC & yes (min gap $1.8\times10^{-2}\,t$ at generic $k$) & yes \\
$\Delta/t$ & 1.5 & 1.5 \\
Max effective spin pol.\ $p_{k\parallel}^{\max}$ & \textbf{0.99997} & ``nearly 100\%'' \\
Fraction of FS cells $p\ge0.90$ (full grid) & 0.144 & broad area \\
Fraction $p\ge0.90$, $E_F\ge-0.5t$ & 0.108 & broad area \\
Fraction $p\ge0.99$ (full grid) & 0.047 & --- \\
\bottomrule
\end{tabular}
\caption{Grid: $n_{k_y}=41$, $n_{k_x}=601$, $n_{E_F}=41$, $E_F\in[-2.2,2.2]\,t$.}
\end{table}

The maximum polarization reaches essentially 100\% ($0.99997$), directly
confirming the qualitative headline that noncollinear AFM conduction channels can
be fully spin polarized. Regions of $p\ge0.90$ concentrate along the
$k_y=0$/near-axis lines away from the zone center---consistent with the paper's
observation (Fig.~1c, Fig.~3e) that polarization is largest away from $\Gamma$ and
reduced near the zone center where multiple noncollinear bands overlap. Our
whole-grid ``broad area'' fraction ($\sim14\%$) is lower than a visual reading of
the paper's map because our coarse Fermi-crossing counter includes the many-band
overlap regions and the full energy window; the physically decisive result---that
$p_{k\parallel}\to1$ is reached at all---is unambiguous.

\section{Verdict}
\textbf{REPLICATED} (tight-binding headline). The from-scratch kagome noncollinear
AFM model reproduces the six spin-split bands and achieves $p_{k\parallel}\approx
100\%$, confirming the mechanism the paper proposes. The DFT Mn$_3$GaN band
structure and the ETMR $\sim10^4\%$ ballistic-tunneling calculation through
SrTiO$_3$ were intentionally not attempted (DFT out of scope for this fast pass).

\section*{Credit}
Physics scaffolding informed by the shared Textures-100 kernels
(\texttt{gobel2024\_sd\_skyrmion\_kubo\_Lz\_kernel.py} s--d lattice pattern and the
kagome-geometry loop-current kernel). Model and analysis written from scratch for
this noncollinear-AFM spin-polarization task.

\end{document}
